\documentclass[11pt,a4paper]{ctexart}
\input{/Users/shaoyikai/Desktop/考研/22408-notes-skill/assets/note-preamble.tex}
\renewcommand{\NoteFooter}{22408 一页纸笔记 · 408 操作系统}
\linespread{1.14}
\begin{document}
\notetitle{文件分配方式选型}{408}{操作系统·文件管理·文件物理结构}
\sechead{一、这道题在考什么}

题干「文件\textbf{数据块}的组织方式」= 文件的\textbf{物理结构}（分配方式），不是逻辑结构。三个关键词定生死：\textbf{CD-ROM} 只读定长、不增不删；\textbf{视频} 量大、要连续码率；\textbf{快速随机播放} 跳过去立刻要能读，寻道最少。

\notebox{\textbf{答案 A 连续结构}：随机访问要「一次算出、直达」，不能多读索引块，更不能顺着链爬。（2013 年 408 第 24 题）官方口径：要保证最短的查询时间，即不能选取链表和索引结构，因此连续结构最优。}

\sechead{二、四种方式的硬指标对比}

访问第 $i$ 个逻辑块的\textbf{总寻道次数}是本题唯一标尺：

\begin{center}\small
\begin{tabularx}{\linewidth}{@{}>{\raggedright\arraybackslash}X>{\raggedright\arraybackslash}X>{\raggedright\arraybackslash}X>{\raggedright\arraybackslash}X>{\raggedright\arraybackslash}X@{}}
\toprule
分配方式 & 定位第 $i$ 块 & 总寻道次数 & 动态增长 & 主要短板 \\
\midrule
连续结构 & 起始块号 $+i$，直接算 & $1$ & 不能，须预留或整体迁移 & 外部碎片、不便增删 \\
链接（隐式） & 从首块沿指针爬 $i$ 步 & $i$ & 能，任意块挂到链尾 & 只支持顺序访问 \\
显式链接（FAT） & 查 FAT 表得块号 & $1\sim2$ & 能 & FAT 需常驻内存 \\
直接索引 & 先读索引块，再读数据块 & $2$ & 能 & 索引块额外占空间 \\
多级索引 & 逐级读 $k$ 级索引块 & $k+1$ & 能 & 小文件也付多级开销 \\
\bottomrule
\end{tabularx}
\end{center}

\notebox{连续结构是唯一「零额外元数据 I/O」的方式：块号由算术得到，不查表、不读索引、不爬链。}

\sechead{三、三步排除}

\begin{enumerate}

\item \textbf{「随机」排除链式}：访问第 $i$ 块要从第 1 块爬 $i$ 步，$i$ 次读盘，播放必卡。

\item \textbf{「快速」排除索引}：直接索引先读索引块，多一次寻道；多级索引再多 $k$ 次。C、D 只是陪衬。

\item \textbf{「CD-ROM」消掉连续的毛病}：连续结构被通用文件系统淘汰，是因为不便增长、增删并有外部碎片；光盘文件只读定长，这个短板不会发生。

\end{enumerate}

\warnline{命题人特意写 CD-ROM，就是在替你删掉连续结构的缺点 —— 这才是题眼，不是「连续最好」这四个字。}

\sechead{四、考场判断流程}

\begin{center}\small
\begin{tabularx}{\linewidth}{@{}>{\raggedright\arraybackslash}X>{\raggedright\arraybackslash}X>{\raggedright\arraybackslash}X@{}}
\toprule
题干信号 & 选谁 & 理由 \\
\midrule
只读 / 一次写入 / 光盘 / 视频点播 / 流式播放 / 要求寻道最少 & 连续结构 & 零额外 I/O，可预读，顺序传输最快 \\
要随机访问，且文件频繁增删改、大小不定 & 索引（UNIX 混合索引） & 随机访问是标配，增长灵活 \\
只要顺序访问，文件频繁增长、几乎不随机 & 链接（隐式 / FAT） & 任意块可挂链，无外部碎片 \\
文件极大，超出索引块寻址范围 & 多级索引 & 用级数换寻址范围 \\
\bottomrule
\end{tabularx}
\end{center}

题干同时出现「随机」与「最快」，基本就是连续结构；追加型日志、磁带备份则连续、链接都行。

\sechead{五、易错点}

\warnline{① \textbf{物理结构 ≠ 逻辑结构}。逻辑结构是顺序文件、索引文件、索引顺序文件、直接文件、散列文件；本题问「数据块组织方式」，答分配方式。}

\warnline{② \textbf{链式分隐式与显式}。隐式链接不支持随机访问（须爬链），FAT 能查表、支持随机访问；做「是否支持随机访问」判断题时务必分开，做本题两者都排不上第一。}

\warnline{③ \textbf{「支持」≠「最好」}。索引结构的随机访问要先付一次索引块寻道，连续结构不付，所以「最快」永远轮不到索引。}

\begin{itemize}

\item 锚点：\textbf{「快」优先选连续，「灵活」优先选索引}；光盘、DVD、蓝光、WORM 上的大文件，教材一律以连续分配为例。

\end{itemize}

\end{document}
